% !TEX program = xelatex
\documentclass[12pt,a4paper]{ctexart}

\usepackage{amsmath,amssymb}
\usepackage{graphicx}
\graphicspath{{../}{../Q1/}{../Q2/}{../Q3/}{../Q4/}{./}}
\usepackage{booktabs}
\usepackage{geometry}
\usepackage{float}
\usepackage{enumitem}
\usepackage[colorlinks=true,linkcolor=black,citecolor=black,urlcolor=black]{hyperref}
\usepackage{listings}
\usepackage{xcolor}

\newfontfamily\codefont{Microsoft YaHei}

\lstset{
    language=Python,
    basicstyle=\codefont\scriptsize,
    numbers=left,
    numberstyle=\tiny,
    frame=single,
    breaklines=true,
    breakatwhitespace=false,
    columns=fullflexible,
    keepspaces=true,
    showstringspaces=false,
    tabsize=4,
    captionpos=t
}
\usepackage{setspace}

\ctexset{
    section/format+={\raggedright},    % 一级标题 \section 左对齐
    subsection/format+={\raggedright},  % 二级小标题 \subsection 左对齐
    subsubsection/format+={\raggedright}% 三级小标题 \subsubsection 左对齐
}

\geometry{
    left=2.5cm,
    right=2.5cm,
    top=2.5cm,
    bottom=2.5cm
}

\setstretch{1.25}

\title{
    \heiti 绿色可回收物回收量动态预测与回收体系稳定性分析
}

\begin{document}

\pagestyle{plain}
\hypersetup{
    pdftitle={绿色可回收物回收量动态预测与回收体系稳定性分析},
    pdfauthor={}
}

\begin{center}
    {\zihao{2}\heiti 绿色可回收物回收量动态预测与回收体系稳定性分析}
\end{center}
\vspace{0.8em}

\begin{abstract}
{\zihao{5}\kaishu
针对上海``沪尚回收''体系回收量动态波动与长效运维管控难题，本文基于9个时间点实测数据，构建Logistic动力学模型，依``正向建模—稳定性判定—灵敏度控制''三层递进求解。基础模型暴露了给定增长率$\gamma=0.022$的系统性高估（$R^2=0.5319$）。引入促进系数$\alpha$与抑制系数$\beta$构建对称拮抗扩展模型，经非线性最小二乘拟合得$r=0.00849\pm0.00066$、$K_{\text{eff}}=8142\pm73$，反解$\alpha=0.016\pm0.009$、$\beta=1.59\pm0.20$（Delta方法），$R^2$提升至0.9915，MAPE降至0.53\%；残差Bootstrap（$B=10^4$）与Delta方法交叉验证了参数估计的稳健性——两种框架的证据方向一致，$\alpha$的显著性处于边界区域（Bootstrap 95\% CI $[0.0006,0.0319]$）。稳定性分析证明物理平衡点$x^*=K_{\text{eff}}$恒为渐近稳定（$\tau\approx118$天），系统在全部有效参数空间内具结构稳定性（$\beta_{\max}=3.40$）。灵敏度分析揭示$K$为量级主导参数（$\pm10\%$偏差使$R^2$骤降至$\sim0.3$），并暴露了``吸引子错位''这一传统稳定性检验无法发现的隐蔽风险。据此提出分层运维方案：精确普查$K$（第一优先）、定向削减$\beta$、促进策略从规模扩张转向质量提升。

{\heiti 关键词}：可回收物回收；Logistic动力学；系统稳定性；分岔分析；参数灵敏度；Bootstrap验证
}
\end{abstract}

\newpage

\section{引言}
在我国“双碳”战略全面推进背景下，生活垃圾分类与可回收物资源化循环利用成为城市生态文明建设的关键抓手。上海打造的“沪尚回收”构建“点—站—场”三级全域回收网络，依托智能回收箱、社区便民点位实现居民就近回收覆盖，但实际运营中回收量受居民参与意愿、宣传推广力度、站点运维水平、清运调度效率等多重因素共同干扰，整体呈现非线性动态波动特征。传统静态规划方法仅能完成存量数据统计，无法刻画回收量长期演化趋势，也难以量化识别系统运行失衡隐患，缺乏支撑长效运营的定量分析工具。
针对上述现实痛点，本文以赛题给出的“沪尚回收”实测时序数据与固定参数为基础\cite{problem}，采用递进式动力学建模思路，依次搭建基础增长模型、引入促进与抑制因子的改进Logistic模型，结合平衡点稳定性理论、单参数扰动灵敏度分析完整解析回收系统演化规律。通过量化各参数作用强度与安全运行边界，区分不同管控措施的收益差异，最终形成分层、可落地的运维优化策略，为国内同类城市绿色回收体系的稳定运营与精细化管理提供定量参考。

本文在参考文献著录核对、文字规范化和程序复现性检查环节使用了
OpenAI 开发的 Codex（GPT-5）作为辅助工具\cite{openai_codex}；模型构建、参数选择、
数学推导、数值结果解释及最终取舍均由参赛队完成并负责。具体使用记录见支撑材料
中的“AI工具使用详情.pdf”。


\section{问题分析}

本文采用递进建模思路完成四项分析：构建基础 Logistic 模型刻画回收自然增长；引入促进、抑制系数建立改进拮抗模型并完成参数拟合；求解平衡点完成系统稳定性边界判定；采用 OAT 单参数扰动开展灵敏度量化，最终输出分层运维优化策略。

\section{模型假设}

为便于建立和分析模型，作如下假设：

\subsection{关于模型结构的基本假设}

\begin{enumerate}
    \item \textbf{Logistic 增长形式假设}：回收量 $x(t)$ 随时间呈 S 型演化，满足一维自治 Logistic 方程 $\dot{x} = f(x)$，反映资源容量上限对增长的约束。

    \item \textbf{连续时间近似假设}：将离散观测数据（9 个时间点）视为连续演化过程，该近似在特征时间尺度 $\tau \approx 118$ 天大于采样间隔的条件下合理。

    \item \textbf{自治系统假设}：$f(x)$ 不显含时间 $t$，即系统动力学规律在研究周期（365 天）内保持不变。
\end{enumerate}

\subsection{关于参数结构与取值的假设}

\begin{enumerate}
    \item[4.] \textbf{促进-抑制解耦假设}：促进系数 $\alpha$ 仅通过 $K_{\text{eff}} = K(1+\alpha)$ 影响容量，抑制系数 $\beta$ 仅通过 $r = \gamma/(1+\beta)$ 影响增长率，二者作用互不交叉。

    \item[5.] \textbf{参数恒定假设}：研究周期内 $\gamma$、$K$、$\alpha$、$\beta$ 均保持恒定，可用常系数 ODE 描述系统演化。

    \item[6.] \textbf{非负性假设}：$\alpha \geq 0$，$\beta \geq 0$，符合促进不降容量、抑制不提速的物理定义。
\end{enumerate}

\subsection{关于数据与空间的假设}

\begin{enumerate}
    \item[7.] \textbf{空间均质化假设}：以全市总回收量 $x(t)$ 为单一状态变量，不区分城区、街道等空间差异。

    \item[8.] \textbf{品类同质化假设}：不区分纸类、塑料、金属、玻璃等品类，统一归总为单一回收量变量。

    \item[9.] \textbf{观测数据代表性假设}：9 个时间点数据能无偏代表系统真实回收水平，观测误差相对模型结构误差可忽略。

    \item[10.] \textbf{参数独立性假设}：问题四的 OAT 灵敏度分析假设各参数波动相互独立，不考虑交互效应（$K$ 与 $\alpha$ 的耦合另行讨论）。
\end{enumerate}

\section{符号说明}

以下符号按其在文中出现的先后顺序排列。注：所有回收量单位统一使用"吨/天"，时间单位统一使用"天"。

\begin{table}[H]
\centering


\begin{tabular}{cclc}
\toprule
符号 & 类型 & 含义 & 单位 \\
\midrule
$t$      & 自变量 & 时间（以竞赛数据起始日为 $t=0$） & 天 \\
$x(t)$   & 状态变量 & 第 $t$ 天的全市可回收物回收量 & 吨/天 \\
$x_0$    & 初始条件 & 初始回收量（$t=0$ 时观测值） & 吨/天 \\
$K$      & 固定参数 & 基准潜在饱和容量（由题目表 (c) 给出） & 吨/天 \\
$\gamma$ & 固定参数 & 固有增长率（由题目表 (c) 给出） & 天$^{-1}$ \\
$\alpha$  & 导出参数 & 促进系数（反映站点扩张、宣传引导的正向效应） & 无量纲 \\
$\beta$   & 导出参数 & 抑制系数（反映清运滞后、运维不足的负向效应） & 无量纲 \\
$r$      & 有效参数 & 有效增长率，$r = \gamma/(1+\beta)$ & 天$^{-1}$ \\
$K_{\text{eff}}$ & 有效参数 & 有效饱和容量，$K_{\text{eff}} = K(1+\alpha)$ & 吨/天 \\
$A$      & 导出常数 & Logistic 解析解简化系数，$A = K/x_0 - 1$ & 无量纲 \\
$f(x)$   & 向量场 & 一维自治系统右端函数，$\dot{x}=f(x)$ & 吨/天$^2$ \\
$\mathcal{D}$ & 定义域 & 系统物理可达状态集，$\mathcal{D}=[x_0, K_{\text{eff}}]$ & 吨/天 \\
$x^*$    & 平衡点 & 满足 $f(x^*)=0$ 的不动点 & 吨/天 \\
$\lambda$ & 特征值 & Lyapunov 线性化特征指数，$\lambda=f'(x^*)$ & 天$^{-1}$ \\
$\tau$    & 特征时间 & 小扰动衰减至 $1/e$ 的松弛时间，$\tau=1/|f'(x^*)|$ & 天 \\
$\eta(t)$ & 扰动量 & 相对平衡点的偏移量，$\eta(t)=x(t)-x^*$ & 吨/天 \\
$r_{\min}$ & 性能阈值 & 运营可接受的最低有效增长率 & 天$^{-1}$ \\
$\varepsilon_{a,b}$ & 弹性 & $b$ 对 $a$ 的弹性系数，$\varepsilon_{a,b} = (\partial a/\partial b) \cdot (b/a)$ & 无量纲 \\
\bottomrule
\end{tabular}
\caption{模型变量与参数一览}
\label{tab:symbols}

\end{table}

\section{模型建立与求解}

\subsection{问题一：基础动态模型}

\subsubsection{模型建立}

考虑到回收体系受到资源容量、人口规模及运营能力等因素限制，回收量通常表现为前期快速增长、后期逐渐趋于饱和的 S 型变化规律。因此，本文采用经典 Logistic 模型描述回收量随时间的动态演化过程。基于经典 Logistic 增长理论\cite{verhulst}，考虑到回收系统在实际运行过程中存在资源
容量约束与增长速率限制，假设回收量 $x(t)$ 随时间呈现先快速增长、
后逐渐趋缓并最终趋于稳定的动态演化规律。其中，固有增长率 $\gamma$
 描述系统在初始阶段的增长潜力，环境容量 $K$ 表示在现有条件下系统能
 够达到的最大稳定回收水平。由于本节主要用于建立基础预测模型，暂不考
 虑外部促进因素和内部抑制机制，仅保留两个核心约束因素：

\begin{equation}\label{eq:ode_q1}
\frac{dx}{dt} = \gamma \cdot x \cdot \left(1 - \frac{x}{K}\right), \qquad x(0) = x_0
\end{equation}

解析解：

\begin{equation}\label{eq:solution}
x(t) = \frac{K}{1 + \left(\dfrac{K}{x_0} - 1\right) e^{-\gamma t}}
     = \frac{K}{1 + A \cdot e^{-\gamma t}}, \qquad A = \frac{K}{x_0} - 1
\end{equation}

\begin{figure}[htbp]
    \centering
    \includegraphics[width=0.75\linewidth]{logistic_fit.png}
    \caption{基础 Logistic 模型拟合曲线}
    \label{fig:logistic_fit}
\end{figure}

如图\ref{fig:logistic_fit}所示，基础模型呈现典型的S型增长曲线，但由于
$\gamma$ 由题目给定而非数据拟合，预测值在早期阶段明显偏高。

\textbf{注意}：由于本模型中的参数 $\gamma$、$K$ 和初始状态 
$x_0$ 均来源于已知条件，因此该模型属于无参数调整的确定性模型。
模型预测结果完全由先验参数决定，不通过数据拟合进一步优化参数，
从而能够用于检验给定理论增长机制与实际恢复过程之间的一致性。

\subsubsection{参数}

\begin{table}[htbp]
\centering

\begin{tabular}{lrll}
\toprule
参数 & 符号 & 数值 & 单位 \\
\midrule
饱和容量      & $K$      & 8011   & 吨/天 \\
固有增长率    & $\gamma$ & 0.022  & 天$^{-1}$ \\
初始回收量    & $x_0$    & 6314   & 吨/天 \\
导出常数      & $A$      & 0.2688 & 无量纲 \\
\bottomrule
\end{tabular}
\caption{问题一 基础 Logistic 模型参数}
\label{tab:q1_params}
\end{table}

\subsubsection{模型评估}

为定量分析基础 Logistic 模型对实际恢复过程的描述能力，
本文选取决定系数 $R^2$、平均绝对误差（MAE）以及平均绝对百分比误差（MAPE）
作为模型评价指标。

其中，$R^2$ 用于衡量模型对数据整体变化趋势的解释能力，
MAE 和 MAPE 分别从绝对误差和相对误差角度评价模型预测结果的偏差程度。

\begin{table}[htbp]
\centering

\begin{tabular}{lc}
\toprule
指标 & 数值 \\
\midrule
$R^2$          & 0.5319 \\
$R_{\text{adj}}^2$ & 0.5319 \\
MAE（吨/天）   & 315.6 \\
RMSE（吨/天）  & 385.4 \\
MAPE          & 4.42\% \\
\bottomrule
\end{tabular}
\caption{问题一 基础 Logistic 模型评估指标}
\label{tab:q1_metrics}
{\small 注：$R_{\text{adj}}^2 = R^2$ 系因模型无拟合参数（$p=0$），由$R_{\text{adj}}^2 = 1 - \frac{n-1}{n-p-1}(1-R^2)$知二者公式等价，非退化。}
\end{table}

\subsubsection{模型结果分析}
\paragraph{预测与残差}

\begin{table}[htbp]
\centering

\begin{tabular}{crrrr}
\toprule
$t$（天） & 实测值（吨/天） & 预测值（吨/天） & 残差（吨/天） & 相对误差 \\
\midrule
0   & 6314 & 6314.0 & $-0.0$   & $-0.00\%$ \\
30  & 6542 & 7033.9 & $-491.9$ & $-7.52\%$ \\
60  & 6875 & 7474.4 & $-599.4$ & $-8.72\%$ \\
90  & 7173 & 7724.4 & $-551.4$ & $-7.69\%$ \\
120 & 7368 & 7860.2 & $-492.2$ & $-6.68\%$ \\
150 & 7591 & 7932.4 & $-341.4$ & $-4.50\%$ \\
180 & 7724 & 7970.2 & $-246.2$ & $-3.19\%$ \\
270 & 7896 & 8005.3 & $-109.3$ & $-1.38\%$ \\
365 & 8002 & 8010.3 & $-8.3$   & $-0.10\%$ \\
\bottomrule
\end{tabular}
\caption{问题一 预测值与残差分析}
\label{tab:q1_predictions}
\end{table}

\textbf{系统性偏差}：从残差分布可以发现，预测值在整个时间区间内均高于实际观测值，
说明基础 Logistic 模型存在一定程度的系统性高估现象。

其中，最大负残差出现在 $t=60$ 天，对应偏差为
$-599.4$ 吨/天。

该结果表明，在恢复初期阶段，实际系统增长速度低于理论模型预设水平。
产生该偏差的原因可能包括资源配置限制、运行效率波动以及外部环境因素影响，
因此固定增长率 $\gamma=0.022$ 难以完全反映真实恢复过程中的动态变化。

\subsection{问题二：引入促进与抑制作用的改进模型}

\subsubsection{模型建立}

为进一步提高基础 Logistic 模型对实际恢复过程的描述能力，
本文在 问题一 模型基础上引入两个具有拮抗作用的调节参数：
促进系数 $\alpha$ 与抑制系数 $\beta$。

其中，$\alpha$ 表征外部促进因素对系统潜在恢复能力的增强作用，
通过提高有效环境容量 $K_{\mathrm{eff}}$ 改变系统长期稳定水平；
$\beta$ 表征内部限制因素对恢复过程的约束作用，
通过降低有效增长率 $r$ 调节系统接近稳定状态的速度。

二者作用方向相反，但共同决定系统动态演化过程，
因此形成具有促进—抑制拮抗机制的扩展 Logistic 模型：

\begin{equation}\label{eq:ode_q2}
\frac{dx}{dt} = \frac{\gamma}{1+\beta} \cdot x \cdot \left(1 - \frac{x}{K(1+\alpha)}\right)
\end{equation}

\begin{itemize}
  \item $\alpha\ge0$ 表示外部促进作用，其通过提高系统潜在承载能力扩大有效环境容量：  $\;\rightarrow\; K_{\text{eff}} = K(1+\alpha)$
  \item $\beta\ge0$ 表示内部抑制作用，其通过降低系统有效增长率影响恢复速度： $\;\rightarrow\; r = \gamma/(1+\beta)$
\end{itemize}

两个调节参数均采用 $(1+c)^{\pm1}$ 的形式作用于模型参数，
从数学结构上保持参数调整的对称性。

其中，$\alpha$ 主要影响系统长期稳定水平，
而 $\beta$ 主要影响系统达到稳定状态的速度。

当 $\alpha = \beta = 0$ 时，扩展模型退化为 问题一 中的基础 Logistic 模型。二者拮抗方向相反：
$\alpha$ 增大饱和水平，$\beta$ 降低趋近速度。

\subsubsection{标准 Logistic 表达式}

经过参数变换后，扩展模型仍保持 Logistic 模型的标准形式，
但其有效环境容量和有效增长率均由调节参数动态决定。

因此，该模型既继承了 Logistic 模型的稳定收敛特性，
又能够描述外部促进因素和内部限制因素共同作用下的恢复过程。

\begin{equation}\label{eq:logistic_q2}
\frac{dx}{dt} = r \cdot x \cdot \left(1 - \frac{x}{K_{\text{eff}}}\right)
\end{equation}
\begin{equation}
r = \frac{\gamma}{1+\beta}, \qquad K_{\text{eff}} = K(1+\alpha)
\end{equation}

解析解：
\begin{equation}\label{eq:solution_q2}
x(t) = \frac{K_{\text{eff}}}{1 + \left(\dfrac{K_{\text{eff}}}{x_0} - 1\right) e^{-r t}}
\end{equation}

\begin{figure}[htbp]
    \centering
    \includegraphics[width=0.75\linewidth]{model_q2.png}
    \caption{引入促进-抑制作用的改进模型拟合曲线}
    \label{fig:model_q2}
\end{figure}

如图\ref{fig:model_q2}所示，引入$\alpha$和$\beta$后拟合曲线与实测数据高度吻合，
残差显著缩小且符号正负交替，无明显系统偏差。

\subsubsection{参数}

由于扩展模型最终可以转化为标准 Logistic 形式，
因此首先对有效增长率 $r$ 和有效环境容量 $K_{\text{eff}}$ 进行参数识别，
再根据参数映射关系反推出 $\alpha$ 和 $\beta$。

\paragraph{拟合方法与算法}

采用\textbf{非线性最小二乘法}（Nonlinear Least Squares, NLS）进行参数估计。
以残差平方和（Residual Sum of Squares, RSS）为目标函数：
\begin{equation}
\min_{r,\,K_{\text{eff}}} \;
\sum_{i=1}^{n} \bigl(y(t_i) - x(t_i;\, r,\,K_{\text{eff}})\bigr)^2,\qquad
n = 9
\end{equation}

根据模型设计要求，当拟合结果满足
$r\le\gamma$ 且 $K_{\text{eff}}\ge K$ 时，
可进一步保证 $\beta\ge0$ 和 $\alpha\ge0$。

该约束条件符合实际恢复过程中的物理意义，
即促进因素不会降低系统容量，而抑制因素不会提高增长速率。


\begin{table}[htbp]
\centering
\caption{问题二 扩展模型参数拟合结果}
\label{tab:q2_params}
\begin{tabular}{llrl}
\toprule
类别 & 符号 & 拟合值 & 单位 \\
\midrule
固定  & $K$                & 8011        & 吨/天 \\
固定  & $\gamma$           & 0.022       & 天$^{-1}$ \\
固定  & $x_0$              & 6314        & 吨/天 \\
\midrule
\textbf{拟合} & $r$                & $\mathbf{0.00849} \pm 0.00066$ & 天$^{-1}$ \\
\textbf{拟合} & $K_{\text{eff}}$   & $\mathbf{8142} \pm 73$ & 吨/天 \\
\midrule
导出  & $\alpha$           & $\mathbf{0.0164} \pm 0.0091$ & 无量纲 \\
导出  & $\beta$            & $\mathbf{1.59} \pm 0.20$ & 无量纲 \\
\bottomrule
\end{tabular}
\end{table}

\textbf{导出参数的不确定度传播（Delta方法）}：$\alpha$ 和 $\beta$ 由拟合参数 $r$、$K_{\text{eff}}$ 经非线性映射导出，其标准误差需通过一阶Taylor展开（Delta方法）传播。由 $\beta = \gamma/r - 1$ 和 $\alpha = K_{\text{eff}}/K - 1$：
\begin{equation}
\sigma_\beta = \left|\frac{\partial\beta}{\partial r}\right| \sigma_r
             = \frac{\gamma}{r^2}\,\sigma_r
             = \frac{0.022}{(0.00849)^2} \times 0.00066
             \approx 0.20
\end{equation}
\begin{equation}
\sigma_\alpha = \left|\frac{\partial\alpha}{\partial K_{\text{eff}}}\right| \sigma_{K_{\text{eff}}}
              = \frac{1}{K}\,\sigma_{K_{\text{eff}}}
              = \frac{1}{8011} \times 73
              \approx 0.0091
\end{equation}
该线性近似在 $r$ 和 $K_{\text{eff}}$ 的相对不确定度较小（$\sigma_r/r \approx 7.8\%$，
$\sigma_{K_{\text{eff}}}/K_{\text{eff}} \approx 0.9\%$）的条件下是合理的。
\begin{equation}
\begin{aligned}
\operatorname{Cov}(\alpha,\beta) &\approx 0,\\
&\text{其中 $\alpha$ 仅依赖 $K_{\text{eff}}$，$\beta$ 仅依赖 $r$；}\\
&\text{$\operatorname{Cov}(r,K_{\text{eff}})$ 应由拟合协方差矩阵给出，此处近似取零。}
\end{aligned}
\end{equation}

\textbf{$\alpha$ 的统计显著性}：小样本下（$n-p=7$）采用 $t$ 分布构造置信区间。
取 $t_{0.025,7}=2.365$，$\alpha$ 的 95\% 置信区间为
$\alpha \pm t_{0.025,7}\,\sigma_\alpha = 0.0164 \pm 2.365 \times 0.0091 = [-0.0051,\;0.0379]$。
该区间包含零，说明在 95\% 置信水平下\textbf{不能排除 $\alpha = 0$ 的原假设}——
即促进系数对有效容量的正向贡献在统计上不显著。
这不意味着促进策略无效（见问题四讨论——$\alpha$ 仅捕获对容量的直接效应，
宣传引导等促进手段可能通过 $\gamma$ 和 $\beta$ 间接发挥作用），
但表明在当前数据精度（$n=9$）下，$\alpha$ 的点估计 $0.0164$ 对应的容量提升（约 $1.6\%$）
不应被过度解读为确凿的系统性增益。$\beta$ 的 95\% 置信区间为
$\beta \pm t_{0.025,7}\,\sigma_\beta = 1.59 \pm 2.365 \times 0.20 = [1.12,\;2.07]$，
显著大于零——抑制效应的存在具有统计显著性。\textbf{需注意}：上述渐近置信区间依赖大样本近似（$n \to \infty$），在 $n=9$ 下其可靠性需经独立交叉验证，不宜直接作为定论。

\textbf{Bootstrap 交叉验证}：残差 Bootstrap\cite{efron} 以经验分布替代未知的真实误差分布来模拟参数的抽样分布，不依赖渐近正态假设——但其自身在 $n=9$ 下也受限于两个事实：（1）仅 9 个残差的离散经验分布对真实误差分布的逼近精度有限；（2）误差 i.i.d. 假设在小样本下无法检验。因此，Bootstrap 在此的定位不是替代 Delta 方法的''更优解''，而是提供一条不依赖渐近假说的\textbf{独立参照线}——若两种基于不同前提的方法给出相近结果，则推断可信度增强；若差异显著，则两种方法均需被质疑。

以原拟合残差为经验分布，有放回重抽样生成 $B=10\,000$ 组伪数据，每组独立重拟合 $r$ 和 $K_{\text{eff}}$，由经验分位数直接读出 95\% 置信区间：
\begin{itemize}
  \item $r$：Bootstrap 95\% CI $[0.00734,\;0.00951]$（渐近 $[0.00694,\;0.01003]$）
  \item $K_{\text{eff}}$：Bootstrap 95\% CI $[8016,\;8267]$（渐近 $[7970,\;8314]$）
  \item $\alpha$：Bootstrap 95\% CI $[0.0006,\;0.0319]$（渐近 $[-0.0051,\;0.0379]$）
  \item $\beta$：Bootstrap 95\% CI $[1.31,\;2.00]$（渐近 $[1.12,\;2.07]$）
\end{itemize}
Bootstrap 与 Delta 方法的 CI 量级一致（SE 比值 0.85--0.88），方向相同——两种基于完全不同前提的方法相互印证，增强了参数推断的可信度。原担忧的``渐近近似偏乐观''并未出现，Delta 方法在实际数据中反而略有保守。

\textbf{$\alpha$ 显著性的交叉验证判断}：渐近 CI 包含零而 Bootstrap CI 排除零，两种方法在 $\alpha$ 的显著性上给出了不同的边界判断。不应简单地以 Bootstrap 为准而否定 Delta——更审慎的解读是：\textbf{两种方法的证据指向$\alpha > 0$的方向，但信号强度处于边缘区域}（Bootstrap CI 下限 $0.0006$ 仅勉强跨过零线）。这一判断的保守程度是恰当的——它既不否认促进效应的存在，也不过度依赖单一统计方法的边界判决。\textbf{统计边缘显著 $\neq$ 实际显著}：无论哪种 CI，$\alpha$ 的弹性均为 $\varepsilon_{K_{\text{eff}},\alpha}=0.016 \ll 1$——促进效应对容量的边际贡献已接近极限这一物理判断，在 Delta 和 Bootstrap 两种框架下完全一致，不受显著性边界分歧的影响。

\textbf{参数物理意义分析}：
\begin{itemize}
  \item $\alpha = 0.0164 \pm 0.0091$：点估计为正，促进效应对有效容量的贡献约为 $+1.64\%$。
  渐近 95\% CI $[-0.0051,\;0.0379]$ 包含零，Bootstrap 95\% CI $[0.0006,\;0.0319]$ 不包含零——
  两种方法的证据方向一致（均指向 $\alpha > 0$），但信号处于统计显著性边界，不宜对 $\alpha$ 的精确数值做过度解读。
  无论哪种框架，弹性 $\varepsilon_{K_{\text{eff}},\alpha}=0.016 \ll 1$ 均表明促进效应对容量的边际贡献已接近极限。
  这不意味着促进策略本身无效（见问题四——促进手段可能通过提升 $\gamma$ 或降低 $\beta$ 间接发挥作用），
  而是表明 $\alpha$ 对容量的直接贡献在当前覆盖密度下已进入收益递减区间。
  \item $\beta = 1.59 \pm 0.20$（渐近 95\% CI $[1.12,\;2.07]$，Bootstrap 95\% CI $[1.31,\;2.00]$，均显著大于零）：
  系统受到统计显著的抑制作用，有效增长率降低至 $r = \gamma/(1+\beta) \approx 0.00849$，
  仅为固有增长率的 $38.6\%$。该结果说明实际恢复过程中的增长速度受到较强限制。
  \item $r/\gamma = 1/(1+\beta) \approx 0.386$：增长率修正幅度（约 $-61.4\%$）
  明显大于容量修正幅度（约 $+1.6\%$），内部限制因素对系统演化过程起主要作用。
\end{itemize}

\subsubsection{计算}

为了验证扩展模型相对于基础 Logistic 模型的改进效果，
本文将 问题一 基础模型与引入 $\alpha,\beta$ 参数后的扩展模型进行对比。

通过比较两类模型的拟合指标，
分析调节参数对于提高预测精度和降低系统误差的有效性。

\begin{table}[htbp]
\centering
\caption{问题一 与 问题二 模型拟合指标对比}
\label{tab:q1_q2_comparison}
\begin{tabular}{lccc}
\toprule
指标 & 问题一（无 $\alpha,\beta$） & 问题二（含 $\alpha,\beta$） & 改进幅度 \\
\midrule
$R^2$         & 0.5319  & \textbf{0.9915} & 提升 $0.4596$ \\
$R_{\text{adj}}^2$ & 0.5319  & \textbf{0.9887} & 提升 $0.4568$ \\
MAE（吨/天）  & 315.6   & \textbf{38.3}   & $-87.9\%$ \\
RMSE（吨/天） & 385.4   & \textbf{51.9}   & $-86.5\%$ \\
MAPE         & 4.42\%  & \textbf{0.53\%}  & $-88.0\%$ \\
\bottomrule
\end{tabular}
\end{table}

从预测结果可以看出，引入 $\alpha$ 和 $\beta$ 后，
模型预测值与实际观测值之间的偏差明显降低。

相比 问题一 中持续存在的系统性高估现象，
扩展模型在整个时间范围内均保持较小残差，
说明调节参数能够有效修正基础 Logistic 模型中增长率设定偏高的问题。

\subsubsection{预测}

\begin{table}[htbp]
\centering
\caption{问题二 预测值与残差分析}
\label{tab:q2_predictions}
\begin{tabular}{crrrr}
\toprule
$t$（天） & 实测值（吨/天） & 预测值（吨/天） & 残差（吨/天） & 相对误差 \\
\midrule
0   & 6314 & 6314.0 & $-0.0$    & $-0.00\%$ \\
30  & 6542 & 6649.5 & $-107.5$  & $-1.64\%$ \\
60  & 6875 & 6935.3 & $-60.3$   & $-0.88\%$ \\
90  & 7173 & 7174.2 & $-1.2$    & $-0.02\%$ \\
120 & 7368 & 7371.2 & $-3.2$    & $-0.04\%$ \\
150 & 7591 & 7531.5 & $+59.5$   & $+0.78\%$ \\
180 & 7724 & 7660.6 & $+63.4$   & $+0.82\%$ \\
270 & 7896 & 7910.5 & $-14.5$   & $-0.18\%$ \\
365 & 8002 & 8037.1 & $-35.1$   & $-0.44\%$ \\
\bottomrule
\end{tabular}
\end{table}

\subsection{问题三：平衡点与稳定性分析}

问题一 与 问题二 完成了模型的建立与参数估计——它们回答了``系统服从何种方程''和``参数取何值能最好地拟合数据''这两个\textbf{正向建模}问题。然而，对于一座城市的回收体系管理者而言，更关键的问题是\textbf{定性层面的}：系统在参数空间内是否始终保持渐近稳定？在何种参数条件下可能发生失稳？安全运行的边界在哪里？$R^2 = 0.9915$ 只能说明模型精确描述了已经观测到的行为——它不能保证系统在未来不会因为参数漂移或外部冲击而失稳。

从本节开始，分析进入动力系统理论的\textbf{动态系统分析}范式。我们将运用相线分析（phase line analysis）判断系统的长期行为模式，借助 Lyapunov 线性化方法给出收敛速率的定量估计，最后通过分岔理论（bifurcation theory）对参数空间进行彻底的结构安全性审视——逐一遍历所有可能的分岔类型，证明每一种在当前物理条件下均不可能发生。

\subsubsection{模型与物理定义域}

问题二 建立的对称拮抗模型可化为标准一维自治 Logistic 系统
$\dot{x} = f(x)$，其向量场为：
\begin{equation}\label{eq:ode_q3}
f(x) = \frac{\gamma}{1+\beta} \cdot x \cdot \left(1 - \frac{x}{K(1+\alpha)}\right)
     = r x \left(1 - \frac{x}{K_{\text{eff}}}\right), \qquad
r = \frac{\gamma}{1+\beta},\;\; K_{\text{eff}} = K(1+\alpha)
\end{equation}
其中 $\alpha \geq 0$, $\beta \geq 0$ 由参数定义保证。
问题二 拟合给出 $r = 0.00849$, $K_{\text{eff}} = 8142$,
$\alpha = 0.0164$, $\beta = 1.59$。

在展开稳定性分析之前，必须明确系统的\textbf{物理可达状态集}。
由初始条件 $x(0) = x_0 = 6314 < K_{\text{eff}} = 8142$，在 logistic 动力学
$\dot{x} = f(x)$ 作用下，$f(x) > 0$ 对任意 $x \in [x_0, K_{\text{eff}})$ 成立，
故 $x(t)$ 从 $x_0$ 出发后单调递增且以 $K_{\text{eff}}$ 为上界——
系统永远不会自发进入 $x > K_{\text{eff}}$ 的区域。定义物理可达状态集为：
\begin{equation}
\mathcal{D} = [x_0,\; K_{\text{eff}}] = [6314,\; 8142]
\end{equation}
（注：$x > K_{\text{eff}}$ 时 $f(x) < 0$，该区域在数学上存在且相线分析中会涉及，
但由上述单调性，从 $x_0$ 出发的物理轨迹永远无法到达。）
\textbf{这一约束对后续稳定性分析具有决定性意义}：数学上两个平衡点中，
只有 $x^*_2 = K_{\text{eff}}$ 落在物理可达集内（$x^*_1 = 0$ 在集外，物理不可达）。

\subsubsection{平衡点求解与物理解释}

令 $f(x^*) = 0$：
\begin{equation}
r x^* \left(1 - \frac{x^*}{K_{\text{eff}}}\right) = 0
\quad\Rightarrow\quad x^*_1 = 0,\;\; x^*_2 = K_{\text{eff}} = 8142
\end{equation}

方程在数学上给出两个平衡点，但两者的物理地位截然不同：
\begin{enumerate}
  \item $x^*_1 = 0$：位于物理可达集 $\mathcal{D}$ 之外（$0 < x_0 = 6314$）。
        回收量从 $x_0$ 出发后单调递增，\textbf{无法}降至 $0$——该平衡点在物理上是不可达的。
  \item $x^*_2 = K_{\text{eff}} = 8142$：位于定义域内，是唯一的物理平衡点。
        代表全市可回收物回收的饱和稳定状态。
\end{enumerate}

\textbf{因此，在物理定义域内，系统有且仅有一个平衡点} $x^* = K_{\text{eff}}$。
后续分析将证明，该平衡点是渐近稳定的，且稳定性在参数扰动下不会丧失。

\subsubsection{相线分析与定性稳定性}

一维自治系统 $\dot{x} = f(x)$ 的动力学行为完全由其向量场符号决定。
逐因子分析 $f(x)$ 的符号（第二区间虽不在物理可达集内，列入以补全数学结构）：

\begin{center}
\begin{tabular}{c|c|c|c|c}
\toprule
区间 & $r$ & $x$ & $1 - x/K_{\text{eff}}$ & $f(x)$ \\
\midrule
$[x_0,\; K_{\text{eff}})$ & $+$ & $+$ & $+$ & $+$（回收量增长） \\
$(K_{\text{eff}},\; +\infty)$ & $+$ & $+$ & $-$ & $-$（回收量衰减） \\
\bottomrule
\end{tabular}
\end{center}

据此画出\textbf{物理相线}（仅显示定义域内部分）：
\begin{center}
$\xrightarrow{\hspace{2.5cm}} \xrightarrow{\hspace{2.5cm}}
\;\bullet\; \xleftarrow{\hspace{2.5cm}} \xleftarrow{\hspace{2.5cm}}$ \\[4pt]
$x_0 = 6314$\hspace{2cm}$K_{\text{eff}} = 8142$\hspace{2cm}
\end{center}

箭头从两侧均指向 $K_{\text{eff}}$，表明：
\begin{itemize}
  \item 当 $x(t) < K_{\text{eff}}$（回收量尚未饱和），$\dot{x} > 0$，回收量单调递增向 $K_{\text{eff}}$
  \item 当 $x(t) > K_{\text{eff}}$（超出饱和极限的理论情形），$\dot{x} < 0$，回收量单调递减回 $K_{\text{eff}}$
\end{itemize}

\textbf{定性结论}：$x^* = K_{\text{eff}}$ 是\textbf{渐近稳定}的——任意物理初值 $x(0) \in \mathcal{D}$
出发的轨迹均单调收敛于该平衡点，\textbf{因此不存在超调、振荡或周期行为}。

这一结论具有深刻的数学基础：对任意一维自治系统 $\dot{x} = f(x)$，相点在
一维实线上只能单向运动，永远不能``回头''，因此周期解在拓扑上不可能存在。
若实际观测数据出现周期性波动（如季节性变化），必由外部驱动力（政策调整、
节假日效应等）所致，而非系统内生动力学。

\subsubsection{线性稳定性分析：定量收敛刻画}

为获得收敛速率的定量估计，对平衡点 $x^* = K_{\text{eff}}$ 执行 Lyapunov
线性化（亦称间接法）。令 $\eta(t) = x(t) - x^*$ 为小扰动量，对 $f(x)$
在 $x^*$ 处做一阶 Taylor 展开：
\begin{equation}
\dot{\eta} = f(x^* + \eta) = f(x^*) + f'(x^*)\eta + O(\eta^2)
           \approx f'(x^*)\,\eta
\end{equation}
其中 $f(x^*) = 0$（平衡点定义），$O(\eta^2)$ 项在小扰动下可忽略。
令 $\lambda = f'(x^*)$，得一阶线性微分方程 $\dot{\eta} = \lambda \eta$，
解为 $\eta(t) = \eta_0 e^{\lambda t}$。

计算模型的导数及其在平衡点处的值：
\begin{equation}
f'(x) = \frac{d}{dx}\left[ r x \left(1 - \frac{x}{K_{\text{eff}}}\right) \right]
       = r\left(1 - \frac{2x}{K_{\text{eff}}}\right)
\end{equation}
\begin{equation}
\begin{aligned}
f'(0) &= r > 0 &&(\text{数学上不稳定，物理不可达}),\\
f'(K_{\text{eff}}) &= r(1 - 2) = -r < 0 &&(\text{渐近稳定}).
\end{aligned}
\end{equation}

$\lambda = f'(K_{\text{eff}}) = -r < 0$ \textbf{恒为负}——物理平衡点具有
指数收敛性质。进一步可定义\textbf{特征时间尺度}：
\begin{equation}
\tau = \frac{1}{|f'(K_{\text{eff}})|} = \frac{1}{r}
     = \frac{1+\beta}{\gamma}
     = \frac{1+1.59}{0.022} \approx 117.7 \;\text{天}
\end{equation}

该时间尺度的物理含义是：小扰动衰减至初始值的 $1/e \approx 37\%$ 约需 118 天。
对于``沪尚回收''体系，这意味着系统的\textbf{自然恢复周期}约为 4 个月——
若受短期运营冲击，系统可在约一个季度的尺度内自然回归正常水平。

\begin{figure}[htbp]
    \centering
    \includegraphics[width=0.75\linewidth]{stability_analysis.png}
    \caption{相线分析、势函数与线性稳定性验证}
    \label{fig:stability_analysis}
\end{figure}

如图\ref{fig:stability_analysis}所示，相线分析（左上）表明$K_{\text{eff}}$两侧向量场均指向平衡点，
势函数（右上）在$K_{\text{eff}}$处取全局极小值，线性化验证（下）确认小扰动呈指数衰减——
三项证据一致支持平衡点的渐近稳定性。

\subsubsection{结构稳定性：参数空间内的分岔分析}

平衡点的稳定性结论回答了``当前参数下系统是否稳定''，但还需回答：\textbf{当参数发生变化时，稳定性是否会丧失？}对此需借助\textbf{分岔理论}\cite{strogatz}审视参数空间的结构安全性。

将模型显式展开：$f(x) = r x - \frac{r}{K_{\text{eff}}} x^2$，无量纲化后得 $\dot{\tilde{x}} = r\tilde{x} - r\tilde{x}^2$，属\textbf{跨临界分岔}标准型族——若 $r$ 可自由穿越零，则 $x^*=0$ 与 $x^*=K_{\text{eff}}$ 的稳定性在 $r=0$ 处交换。然而物理约束锁死了穿越的可能：
\begin{equation}
\beta \geq 0 \;\Rightarrow\; r = \frac{\gamma}{1+\beta} \in (0,\; \gamma],\qquad \gamma = 0.022 > 0
\end{equation}
$r > 0$ \textbf{恒成立}，系统被不可逆地约束于分岔图 $r > 0$ 半平面——$x^* = 0$ 不稳定、$x^* = K_{\text{eff}}$ 稳定的格局不会改变。其余分岔类型亦不适用：鞍-结分岔要求不动点成对湮灭，而 $x^*=0$ 由因式分解产生、不可能消失；叉形分岔要求 $f(-x)=-f(x)$ 的奇对称性，而模型中 $x^2$ 项的存在直接排除了这一可能。

\paragraph{结构稳定性结论}

\begin{center}
\fbox{%
\begin{minipage}{0.88\textwidth}
\centering
\textbf{系统在全部有效参数空间内具有结构稳定性。}\\[4pt]
对任意 $\alpha \geq 0$, $\beta \geq 0$ 和任意物理初值 $x(0) \geq x_0$：\\[2pt]
(1) 物理定义域内恒存在唯一正平衡点 $x^* = K_{\text{eff}} > 0$；\\[2pt]
(2) 该平衡点恒为渐近稳定（$\lambda = -r < 0$）；\\[2pt]
(3) 有效参数空间内不存在任何分岔点——不动点的数量和稳定性类型均不变。\\[4pt]
\textbf{``沪尚回收''体系在数学上不存在分岔失稳——不会发生定性的动力学突变。}
\end{minipage}}
\end{center}

\subsubsection{实用安全边界}

虽然系统具有数学上的结构稳定性，但运行效率要求参数保持在合理范围内。
此处的安全边界由\textbf{性能约束}定义——确保回收体系的收敛速度不低于
运营可接受的最低水平，而非担心数学分岔。

\paragraph{增长率约束}

有效增长率过低将导致系统恢复过慢，影响实际运营效率：
\begin{equation}
r = \frac{\gamma}{1+\beta} \geq r_{\min} \quad\Rightarrow\quad
\beta \leq \frac{\gamma}{r_{\min}} - 1
\end{equation}
$r_{\min}$ 的取值依据为：Logistic 系统恢复至 $99\%$ 饱和水平所需时间满足保守上界
$t_{99\%} \leq -\ln(0.01)/r = \ln(100)/r$（忽略初值因子
$K_{\text{eff}}/x_0-1 \approx 0.29$ 的加速效应，取最不利情形）。
要求收敛时间不超过 $2.5$ 个观测周期（$2.5 \times 365 = 912.5$ 天），
即 $r \geq \ln(100)/912.5 \approx 0.00505$，取整得 $r_{\min} = 0.005$。
代回约束不等式得 $\beta_{\max} = 0.022/0.005 - 1 = 3.40$。

\paragraph{容量约束}
\begin{equation}
K_{\text{eff}} = K(1+\alpha) \geq K_{\min} \quad\Rightarrow\quad
\alpha \geq \frac{K_{\min}}{K} - 1
\end{equation}
取 $K_{\min} = K$（至少不低于基准容量）：$\alpha \geq 0$——由定义自然满足。

\paragraph{安全运行范围与裕度}

\begin{table}[htbp]
\centering
\caption{安全运行范围与裕度指标}
\label{tab:safety_margins}
\begin{tabular}{lcccc}
\toprule
指标 & 下界 & 上界 & 当前值 & 裕度 \\
\midrule
$\alpha$（促进） & $0$（定义） & — & $0.0164$ & — \\
$\beta$（抑制） & $0$（定义） & $3.40$ & $1.59$ & $1.81$ \\
\midrule
$r$ 高于 $r_{\min}$ & \multicolumn{2}{c}{$r_{\min}=0.005$} & $0.00849$ & $0.0035$ \\
$r/\gamma$（有效增长率占比）& \multicolumn{2}{c}{—} & $38.6\%$ & — \\
特征时间尺度 $\tau$ & \multicolumn{2}{c}{—} & $\approx 118$ 天 & — \\
\bottomrule
\end{tabular}
\end{table}

当前 $\beta$ 距性能约束上限仍有 $1.81$ 的充裕裕度（需在当前基础上再增大
$+114\%$ 才触及 $r_{\min}$ 底线），系统在\textbf{收敛速度层面的安全性非常高}。
真正的运营风险不在速度不足，而在 问题四 将揭示的\textbf{容量误判}风险。

\subsubsection{参数解耦与独立调控}

模型结构的另一个重要特性：$\alpha$ 和 $\beta$ \textbf{作用完全解耦}——
本质上是因为促进效应仅作用于有效容量 $K_{\text{eff}}$、抑制效应仅作用于
有效增长率 $r$，二者不交叉。这在分岔理论中是一个重要特性：不存在参数间的
交互分岔（codimension-2 bifurcation），简化了控制策略设计。
\begin{itemize}
  \item $\alpha \to K_{\text{eff}}$（容量维度）：调节促进力度改变饱和容量，
        不影响收敛速度——对应``建站点、扩覆盖''等物理扩张策略
  \item $\beta \to r$（速度维度）：控制抑制力度改变收敛速度，
        不影响饱和容量——对应``提升清运效率、减少运维阻力''等运营优化策略
\end{itemize}

这一解耦特性使得决策者可\textbf{在两个独立维度上分别施策}，
无需担心调控一个参数时意外触发另一参数的不利连锁反应。

\subsection{问题四：参数灵敏度分析}

\subsubsection{问题背景与分析目标}

问题二 针对上海``沪尚回收''体系建立了对称拮抗 Logistic 模型，
引入促进系数 $\alpha$ 和抑制系数 $\beta$ 分别调制有效容量 $K_{\text{eff}} = K(1+\alpha)$
与有效增长率 $r = \gamma/(1+\beta)$，
拟合得到 $R^2 = 0.9915$，较 问题一（$R^2 = 0.5319$）提升显著。
问题三 进一步证明了该系统在数学上具有\textbf{结构稳定性}——
只要 $\alpha, \beta \geq 0$（参数定义本身），
回收量 $x(t)$ 恒趋向正平衡值 $K_{\text{eff}}$，不存在分岔失稳。

然而，从动力系统理论的视角审视，结构稳定性仅保证了系统的\textbf{定性鲁棒性}——
向量场 $f(x)$ 在受到小扰动后相图的拓扑结构不变。但这一概念并未回答
\textbf{定量}问题：吸引子位置（$K_{\text{eff}}$）和收敛速率（$r$）对参数
波动的敏感程度如何？在分岔理论中，当系统远离分岔点时，动力学行为对参数的
变化通常是渐进的、可控的；但当系统接近分岔点时，即使是微小的参数变化
也可能触发定性的跃变。因此，量化系统距分岔点的``距离''以及各参数对系统
输出的弹性，是判断运营安全性的关键补充。基于此，问题四 具体回答以下问题：
\begin{itemize}
  \item 若参数估计存在偏差，模型预测会偏离多远？
  \item 哪些参数的波动对回收量影响最剧烈？
  \item 系统距最近分岔点的``安全距离''有多大？
  \item 在运营资源有限的前提下，应优先精确测量哪些参数、重点调控哪些环节？
\end{itemize}

本问题（问题四）正是为了填补这一空白。
基于表 (c) 规定的 $\pm 10\%$ 单参数扰动（One-At-a-Time, OAT）\cite{saltelli}，
对模型四个基础参数——固有增长率 $\gamma$、抑制系数 $\beta$、促进系数 $\alpha$、
基准容量 $K$——进行局部灵敏度分析，
量化各参数对拟合质量（$R^2$, MAE）和有效参数（$r$, $K_{\text{eff}}$）的影响程度，
为``沪尚回收''体系的长效运维策略提供定量决策依据。

\subsubsection{方法}

\paragraph{模型回顾}

\begin{equation}\label{eq:ode_q4}
\frac{dx}{dt} = \frac{\gamma}{1+\beta} \cdot x \cdot \left(1 - \frac{x}{K(1+\alpha)}\right)
\quad\Rightarrow\quad
\frac{dx}{dt} = r \cdot x \cdot \left(1 - \frac{x}{K_{\text{eff}}}\right)
\end{equation}

解析解：
\begin{equation}\label{eq:solution_q4}
x(t) = \frac{K_{\text{eff}}}{1 + \left(\dfrac{K_{\text{eff}}}{x_0} - 1\right) e^{-r t}},\qquad
r = \frac{\gamma}{1+\beta},\quad K_{\text{eff}} = K(1+\alpha)
\end{equation}

\paragraph{参数物理意义（结合``沪尚回收''实际场景）}

\begin{table}[htbp]
\centering
\caption{参数物理含义（结合``沪尚回收''实际场景）}
\label{tab:param_meaning}
\begin{tabular}{lclp{7cm}}
\toprule
参数 & 符号 & 基准值 & 物理含义 \\
\midrule
固有增长率 & $\gamma$ & 0.022 /天 &
  无外部作用下回收量的自然增长率，反映居民基础参与度
  与政策例行推动效率；对应背景中``居民回收习惯养成周期长''的基线水平 \\

抑制系数 & $\beta$ & 1.59 &
  反映清运延迟、站点运维不足、智能回收箱满溢后未及时清运、
  老年群体操作困难等负面效应的综合强度；$\beta$ 越大则有效增长率越低 \\

促进系数 & $\alpha$ & 0.0164 &
  反映站点覆盖率提升、投递便捷性改善、宣传引导力度加大等
  正向激励的强度；$\alpha$ 越大则有效饱和容量越高 \\

基准容量 & $K$ & 8011 吨/天 &
  全市可回收物潜在饱和上限，由常住人口规模、人均可回收物
  产生量、品类结构等外生因素决定 \\
\bottomrule
\end{tabular}
\end{table}

\paragraph{OAT 扰动方案}

对每个参数 $p \in \{\gamma, \beta, \alpha, K\}$，
分别令 $p \to p \times (1 \pm 0.10)$，其余参数保持基准值不变，
重新计算 $r$, $K_{\text{eff}}$ 及所有拟合指标。

\textbf{基准参数}（来自 问题二 拟合）：

\begin{table}[htbp]
\centering
\caption{问题四 灵敏度分析基准参数}
\label{tab:baseline_params}
\begin{tabular}{lclc}
\toprule
参数 & 符号 & 基准值 & 单位 \\
\midrule
固有增长率 & $\gamma$ & 0.022   & 天$^{-1}$ \\
抑制系数   & $\beta$  & 1.59  & 无量纲 \\
促进系数   & $\alpha$ & 0.0164  & 无量纲 \\
基准容量   & $K$      & 8011    & 吨/天 \\
\midrule
有效增长率 & $r = \gamma/(1+\beta)$ & 0.00849 & 天$^{-1}$ \\
有效容量   & $K_{\text{eff}} = K(1+\alpha)$ & 8142 & 吨/天 \\
\bottomrule
\end{tabular}
\end{table}

基准评估：$R^2 = 0.9915$, MAE $= 38.3$ 吨/天, RMSE $= 51.9$ 吨/天。

\subsubsection{扰动结果与数据解读}

以下对每个参数分别报告 $\pm 10\%$ 扰动结果，并结合``沪尚回收''实际运营场景进行解读。

\paragraph{$\gamma$（固有增长率）$\pm 10\%$}

$\gamma$ 仅通过 $r = \gamma/(1+\beta)$ 影响模型；$K_{\text{eff}}$ 不变。
弹性 $\varepsilon_{r,\gamma} = +1.0000$（线性传递）。

\begin{table}[htbp]
\centering
\caption{$\gamma$（固有增长率）$\pm 10\%$ 扰动对拟合指标的影响}
\label{tab:sens_gamma}
\begin{tabular}{lccc}
\toprule
输出 & $-10\%$ ($\gamma = 0.0198$) & 基准 & $+10\%$ ($\gamma = 0.0242$) \\
\midrule
$r$      & 0.007636 & 0.00849 & 0.009333 \\
$K_{\text{eff}}$ & 8142 & 8142 & 8142 \\
$R^2$    & 0.9810 & 0.9915 & 0.9830 \\
MAE      & 59.8 & 38.3 & 57.7 \\
RMSE     & 77.7 & 51.9 & 73.5 \\
\bottomrule
\end{tabular}
\end{table}

\textbf{数据解读}：
\begin{itemize}
  \item $R^2$ 下降约 0.01（从 0.9915 到 $\sim$0.982），MAE 增加约 20 吨/天。
        变化幅度中等，模型对 $\gamma$ 的估计偏差具有一定容忍度。
  \item 物理原因：数据已处于近饱和区（$x(0)/K_{\text{eff}} = 77.5\%$，$x(365)/K_{\text{eff}} \approx 98\%$），
        增长速率 $r$ 的变化主要影响中期（$t = 30$--$180$ 天）预测值，
        而不改变饱和平台。随着 $t \to \infty$，$\exp(-rt) \to 0$，
        $r$ 对 $x(t)$ 的影响指数衰减——logistic 在近饱和区的``自压缩''效应抑制了 $\gamma$ 扰动的放大。
  \item 实际含义：\textbf{居民基础参与度的短期波动（如季节性变化）不会显著改变长期回收趋势}。
        即使固有人均回收意愿有 $\pm 10\%$ 的估计偏差，模型仍保持 $R^2 > 0.98$。
\end{itemize}

\paragraph{$\beta$（抑制系数）$\pm 10\%$}

$\beta$ 仅通过 $r = \gamma/(1+\beta)$ 影响模型。弹性 $\varepsilon_{r,\beta} = -\beta/(1+\beta) = -0.614$
（非线性衰减——$\beta$ 越大，单位增量对 $r$ 的边际抑制越弱）。

\begin{table}[htbp]
\centering
\caption{$\beta$（抑制系数）$\pm 10\%$ 扰动对拟合指标的影响}
\label{tab:sens_beta}
\begin{tabular}{lccc}
\toprule
输出 & $-10\%$ ($\beta = 1.4336$) & 基准 & $+10\%$ ($\beta = 1.7522$) \\
\midrule
$r$      & 0.009040 & 0.00849 & 0.007994 \\
$K_{\text{eff}}$ & 8142 & 8142 & 8142 \\
$R^2$    & 0.9877 & 0.9915 & 0.9882 \\
MAE      & 49.7 & 38.3 & 48.3 \\
RMSE     & 62.4 & 51.9 & 61.3 \\
\bottomrule
\end{tabular}
\end{table}

\textbf{数据解读}：
\begin{itemize}
  \item $\beta$ 的扰动效果与 $\gamma$ 相近（$R^2$ 降幅 $\sim$0.003--0.004），
        但方向相反——$\beta$ 减小（抑制减弱）使 $r$ 增大，拟合反而略微改善。
        这与 问题二 的结论一致：当前 $\beta = 1.59$ 较大，抑制效应已占主导（$r$ 仅为 $\gamma$ 的 $38.6\%$）。
  \item 弹性 $|\varepsilon_{r,\beta}| = 0.61 < 1$ 意味着 $\beta$ 变化 $1\%$ 仅引起 $r$ 反向变化 $0.61\%$——
        \textbf{抑制效应存在边际递减}：当 $\beta$ 已经较大时，进一步增加抑制的额外伤害有限；
        反之，降低 $\beta$ 的边际收益也递减。
  \item 实际含义：题目背景中描述的``智能回收箱满溢后清运不及时''和``老年群体操作困难''
        这些抑制因素\textbf{已经将有效增长率压至固有水平的 39\%}。
        改善运维效率是有效的，但不宜期待单点突破——消除 $10\%$ 的抑制仅能恢复约 $6.5\%$ 的增长速率。
\end{itemize}

\paragraph{$\alpha$（促进系数）$\pm 10\%$}

$\alpha$ 仅通过 $K_{\text{eff}} = K(1+\alpha)$ 影响模型。
弹性 $\varepsilon_{K_{\text{eff}},\alpha} = \alpha/(1+\alpha) = +0.0161$（极不敏感）。

\begin{table}[htbp]
\centering
\caption{$\alpha$（促进系数）$\pm 10\%$ 扰动对拟合指标的影响}
\label{tab:sens_alpha}
\begin{tabular}{lccc}
\toprule
输出 & $-10\%$ ($\alpha = 0.0147$) & 基准 & $+10\%$ ($\alpha = 0.0180$) \\
\midrule
$r$      & 0.00849 & 0.00849 & 0.00849 \\
$K_{\text{eff}}$ & 8129.1 & 8142 & 8155.3 \\
$R^2$    & 0.9913 & 0.9915 & 0.9913 \\
MAE      & 37.4 & 38.3 & 41.0 \\
RMSE     & 52.4 & 51.9 & 52.4 \\
\bottomrule
\end{tabular}
\end{table}

\textbf{数据解读}：
\begin{itemize}
  \item \textbf{$\alpha$ 的效应近乎为零}：$\pm 10\%$ 扰动下，$R^2$ 变化不超过 $0.0002$，
        $K_{\text{eff}}$ 仅变化 $\pm 13$ 吨/天（$\sim 0.16\%$）。
  \item 数学根源：$\alpha = 0.0164 \ll 1$，$1+\alpha \approx 1.016$，
        促进效应对容量的贡献仅为 $1.64\%$。即使将促进力度翻倍（$+100\%$），
        $K_{\text{eff}}$ 也仅增加约 $1.6\%$。
  \item \textbf{这不是说促进策略无效}，而是说：
        \begin{enumerate}
          \item 当前的站点覆盖率、投递便捷性已经接近其\textbf{对容量的边际贡献极限}——
                全市 1.5 万个服务点 + 4300 台智能回收箱 + 800 个惠民服务点已基本实现社区全覆盖，
                继续增加站点密度对总回收量的提升非常有限（$\alpha$ 对 $K_{\text{eff}}$ 的弹性仅 0.016）；
          \item 然而，促进策略可能通过\textbf{其他渠道}发挥作用——例如提高居民参与度
                （可能改变 $\gamma$，而非 $\alpha$）或降低抑制（减少 $\beta$）。
                模型中 $\alpha$ 仅作用于容量，这不意味着宣传引导的全部价值都被 $\alpha$ 捕获。
        \end{enumerate}
  \item 实际含义：\textbf{以扩大站点数量来提升总回收容量已进入收益递减区间}。
        当前阶段的策略重心应从``广铺网络''转向``精耕运营''——
        例如提高现有站点的清运效率（降低 $\beta$）和改善用户体验（降低 $\beta$），
        而非继续增加站点数量。
\end{itemize}

\textbf{关于 $\alpha$ 结论与 $K$ 的关系}：$\alpha = K_{\text{eff}}/K - 1$ 并非从
数据中独立拟合，而是由 $K_{\text{eff}} \approx 8142$（数据硬约束——晚期数据点
密集分布于饱和平台附近，将 $K_{\text{eff}}$ 锁定于窄区间）与 $K = 8011$
（题目给定值）之差反解得出。因此 $\alpha = 0.016$ 及弹性 $0.016$ 的结论以
$K = 8011$ 的准确性为前提。这恰是下文 $K$ 灵敏度分析将 $K$ 核定列为第一优先级
的根本原因：$K$ 的验证是所有衍生参数结论的共同前提——在 $K$ 经核定确认后，
$\alpha$ 的不敏感属性具有坚实的数据与物理支撑，不构成自指矛盾。

\paragraph{$K$（基准容量）$\pm 10\%$}

$K$ 仅通过 $K_{\text{eff}} = K(1+\alpha)$ 影响模型。
弹性 $\varepsilon_{K_{\text{eff}},K} = +1.0000$（线性——$K$ 变化 $1\% \to K_{\text{eff}}$ 同向变化 $1\%$）。

\begin{table}[htbp]
\centering
\caption{$K$（基准容量）$\pm 10\%$ 扰动对拟合指标的影响}
\label{tab:sens_K}
\begin{tabular}{lccc}
\toprule
输出 & $-10\%$ ($K = 7210$) & 基准 & $+10\%$ ($K = 8812$) \\
\midrule
$r$      & 0.00849 & 0.00849 & 0.00849 \\
$K_{\text{eff}}$ & 7328.0 & 8142 & 8956.4 \\
$R^2$    & 0.2630 & 0.9915 & 0.3256 \\
MAE      & 405.3 & 38.3 & 403.7 \\
RMSE     & 483.6 & 51.9 & 462.6 \\
\bottomrule
\end{tabular}
\end{table}

\textbf{数据解读}：
\begin{itemize}
  \item \textbf{$K$ 是决定性主导参数}：$K \pm 10\%$ 导致 $R^2$ 从 $0.9915$ 骤降至 $0.26$--$0.33$，
        MAE 从 $38.3$ 急剧增大至 $\sim 405$ 吨/天（$\times 10.6$）。
        相比之下，$\gamma, \beta, \alpha$ 的扰动对 $R^2$ 的影响仅为 $K$ 的 $1/30$--$1/170$。
  \item 数学根源：观察数据——$x(365) = 8002$ 已非常接近 $K_{\text{eff}} = 8142$（占 $98.3\%$）。
        所有晚期数据点都密集分布在饱和平台附近。
        $K_{\text{eff}}$ 的任何偏差都会直接转化为饱和预测值的系统性偏差，
        从而对所有 $t \geq 90$ 天的点产生同向残差，导致 $R^2$ 剧烈下降。
        这与 $\gamma$ 或 $\beta$ 的扰动行为截然不同——
        后者只改变曲线的``弯曲程度''而非最终高度，饱和预测值始终收敛于同一个 $K_{\text{eff}}$。

  \item \textbf{$K$ 正负扰动的不对称性}：$K -10\%$ 使 $R^2$ 降至 $0.2630$，
        $K +10\%$ 降至 $0.3256$，降幅差约 $0.06$。
        $K -10\%$ 时有效容量低估，全部预测值均低于实测值（残差全部同号，
        无抵消效应），$R^2$ 下降最为剧烈；$K +10\%$ 时，早期点的负残差
        （预测高出实测）与晚期点的正残差（预测低于实测）方向相反、
        部分互抵，$R^2$ 下降幅度稍弱。
        这一不对称现象并非偶然——其根源与 $K$ 主导性相同：
        所有实测点 $x(t)$ 均位于 $K_{\text{eff}}$ 下方（$x(t) < K_{\text{eff}}$，$\forall t$），
        数据点在 $K_{\text{eff}}$ 两侧的分布本身就高度不对称。
  \item 实际含义：\textbf{基准容量 $K$ 是模型预测有效性的前提条件}。
        $\pm 10\%$ 压力测试表明，若 $K$ 的取值存在量级偏差，所有拟合指标将全面失效——
        $K$ 的准确性是所有衍生参数结论（包括上节 $\alpha$ 的不敏感属性）和策略建议的
        共同前提。题目给定 $K = 8011$ 吨/天为建模提供了基准输入；在实际运营中，
        该值需通过全市可回收物潜力普查加以核定，以确保上述推论链条的基础牢固。
\end{itemize}

\subsubsection{灵敏度排序与参数等级}

\paragraph{对 $R^2$ 的影响}

\begin{table}[htbp]
\centering
\caption{各参数扰动对 $R^2$ 的影响排序（按偏离基线的最大降幅）}
\label{tab:sensitivity_r2}
\begin{tabular}{lcccc}
\toprule
参数 & 基线 $R^2$ & 最差 $R^2$ & $R^2$ 降幅 & 敏感等级 \\
\midrule
$K$      & 0.9915 & 0.2630 & 0.7285 & \textbf{极高}（偏离基线约70个百分点）\\
$\gamma$ & 0.9915 & 0.9810 & 0.0105 & 中等（偏离基线约1个百分点）\\
$\beta$  & 0.9915 & 0.9877 & 0.0038 & 低（偏离基线约0.4个百分点）\\
$\alpha$ & 0.9915 & 0.9913 & 0.0002 & 可忽略 \\
\bottomrule
\end{tabular}
\end{table}

\paragraph{对 MAE 的影响}

\begin{table}[htbp]
\centering
\caption{各参数扰动对 MAE 的影响排序（按偏离基线的最大增幅）}
\label{tab:sensitivity_mae}
\begin{tabular}{lccc}
\toprule
参数 & 基线 MAE（吨/天） & 最差 MAE（吨/天） & MAE 增幅（吨/天） \\
\midrule
$K$      & 38.3 & 405.3 & $+367.0$（$\times 10.6$）\\
$\gamma$ & 38.3 & 59.8  & $+21.5$ \\
$\beta$  & 38.3 & 49.7  & $+11.4$ \\
$\alpha$ & 38.3 & 41.0  & $+2.7$ \\
\bottomrule
\end{tabular}
\end{table}

\textbf{注}：由表\ref{tab:sensitivity_r2}和表\ref{tab:sensitivity_mae}可见，
$K$ 的主导地位是量级上的——$\pm 10\%$ 扰动使 MAE 从 38.3 急剧增大至 $\sim 405$ 吨/天（$\times 10.6$），
而 $\gamma$、$\beta$、$\alpha$ 扰动后 MAE 仍在 37--60 之间。
$K$ 的 $R^2$ 降幅（0.7285）是次敏感参数 $\gamma$（0.0105）的 $\sim 69$ 倍。
所有其余参数的扰动效应在 $K$ 面前均可忽略。

\paragraph{综合敏感等级}

\begin{table}[htbp]
\centering
\caption{参数敏感度等级划分表}
\label{tab:sensitivity_rank}
\begin{tabular}{p{2.2cm} c p{4.6cm} p{5.5cm}}
\toprule
敏感度等级 & 参数 & 弹性特征 & 运营含义 \\
\midrule
\textbf{极度敏感} & $K$ & $\varepsilon(K_{\text{eff}}, K)=1$，数据近饱和 & 必须精确测定；任何估计偏差都将导致模型失效 \\
\textbf{中度敏感} & $\gamma$ & $\varepsilon(r, \gamma)=1$，Logistic饱和区衰减 & 可容忍$\pm10\%$偏差；无需投入大量资源精确估计 \\
\textbf{中度敏感} & $\beta$ & $\varepsilon(r, \beta)=-0.61$，边际递减 & 降低抑制有效但收益递减；应优先消除最严重的瓶颈 \\
\textbf{不敏感} & $\alpha$ & $\varepsilon(K_{\text{eff}}, \alpha)=0.016 \ll 1$ & 促进效应对容量的贡献已饱和；需重新审视促进策略的着力点 \\
\bottomrule
\end{tabular}
\end{table}

\subsubsection{解析弹性汇总}

\begin{table}[htbp]
\centering
\caption{有效参数对基础参数的解析弹性汇总}
\label{tab:elasticity_summary}
\begin{tabular}{lccc}
\toprule
输出 & 输入 & 弹性 $\varepsilon$ & 传递特性 \\
\midrule
$r$ & $\gamma$ & $+1.0000$ & 线性——$\gamma$ 变化 $1\% \to r$ 同向变化 $1\%$ \\
$r$ & $\beta$  & $-0.6143$ & 非线性衰减——$\beta$ 越大边际效应越弱 \\
$K_{\text{eff}}$ & $\alpha$ & $+0.0161$ & 近零——$\alpha \ll 1$ 导致弹性锁死在极小值 \\
$K_{\text{eff}}$ & $K$      & $+1.0000$ & 线性——$K$ 变化 $1\% \to K_{\text{eff}}$ 同向变化 $1\%$ \\
\bottomrule
\end{tabular}
\end{table}

\textbf{弹性解释}：
\begin{itemize}
  \item $|\varepsilon| \approx 0$：输出对该参数不敏感——估计偏差或调控努力的回报极低
  \item $|\varepsilon| = 1$：输出与该参数成比例变化——估计偏差直接传导
  \item $|\varepsilon| > 1$：输出放大参数变化——需格外警惕估计误差
\end{itemize}

值得注意的是，所有解析弹性的绝对值均 $\leq 1$——模型不存在``放大性参数''。
即使是敏感度最高的 $K$，其弹性也仅为 $1$（比例传导，非放大）。
这意味着\textbf{模型在局部线性化意义下是良态的}——不存在微扰导致巨变的分岔或奇点。

\subsubsection{预测包络线分析}

图\ref{fig:sensitivity_tornado}和图\ref{fig:sensitivity_envelope}以两种可视化方法展示了
灵敏度分析的完整结果:

\begin{figure}[htbp]
    \centering
    \includegraphics[width=0.75\linewidth]{sensitivity_tornado.png}
    \caption{龙卷风图：各参数 $\pm 10\%$ 对 $R^2$、MAE、$r$、$K_{\text{eff}}$ 的影响。
    (a)(b) 子图中 $K$ 的影响远超坐标轴范围（以橙色文本框标注离轴数值），意在让读者关注非 $K$ 参数之间的差异；
    (c) $r$ 仅受 $\gamma$ 和 $\beta$ 影响，(d) $K_{\text{eff}}$ 仅受 $\alpha$ 和 $K$ 影响——
    验证了模型的结构解耦特性。}
    \label{fig:sensitivity_tornado}
\end{figure}
\begin{figure}[htbp]
    \centering
    \includegraphics[width=0.75\linewidth]{sensitivity_envelope.png}
    \caption{(a) 所有参数 $\pm 10\%$ 的总预测包络线——包络宽度完全由 $K$ 主导；
    (b) 各参数独立扰动带——$K$（橙色）的带宽远超其余参数，$\alpha$（绿色）几乎不可见；
    (c) 拟合指标的数值弹性——$K$ 对 $R^2$ 的弹性（$\sim 96$）比其余参数高出约两个数量级。}
    \label{fig:sensitivity_envelope}
\end{figure}

\subsubsection{OAT 方法的局限性与补充讨论}

OAT（One-At-a-Time）方法简洁直观、计算成本低，适合作为参数筛选的第一步，但存在三点局限：

\begin{enumerate}
    \item \textbf{忽略参数交互效应}：$K$ 与 $\alpha$ 通过 $K_{\text{eff}} = K(1+\alpha)$ 耦合，OAT 仅逐一扰动单参数。当前 $\alpha \approx 0.016 \ll 1$，交互项可忽略，但此结论依赖具体数据而非方法本身。

    \item \textbf{排序依赖基准点}：OAT 揭示的是当前参数附近的局部敏感性，若基准参数发生变化（如模型迁移至其他城市），排序结果需重新评估。

    \item \textbf{未捕获高阶结构}：弹性系数仅为一阶导数，无法区分线性传导与非线性放大。本文已通过 $\pm10\%$ 数值实验验证 $K$ 的弹性确为线性传导。
\end{enumerate}

后续可采用 Sobol' 全局灵敏度分析\cite{saltelli}量化参数交互效应，或引入贝叶斯方法为 $K$ 赋予先验分布，生成带置信区间的预测曲线。

\subsubsection{与 问题三 稳定性结论的联合分析}

问题三 从动力系统的分岔理论出发，证明了当前参数空间 $(\alpha \geq 0, \beta \geq 0)$
内系统具有结构稳定性：唯一物理平衡点 $x^* = K_{\text{eff}}$ 恒为渐近稳定，
且 $r = \gamma/(1+\beta) > 0$ 恒成立——系统被不可逆地约束于跨临界分岔点
$r = 0$ 的右侧，不存在分岔失稳的任何可能。

然而，\textbf{结构稳定 $\neq$ 运营安全}。系统远离分岔点固然意味着不会发生
定性的动力学突变，但同时也意味着\textbf{另一种风险被掩盖了}：当吸引子
（平衡点 $K_{\text{eff}}$）的位置本身不确定时，系统会以可靠的收敛性
稳定地趋向一个错误的目标值——这种失败模式比突然失稳更加隐蔽，因为
所有的稳定性检验都会通过。问题四 的灵敏度分析正是为了暴露这种\textbf{``稳定性陷阱''}。

联合 问题三 与 问题四 的全部量化证据，可以得出三个核心洞察：

\begin{enumerate}
  \item \textbf{稳定性保证了方向，灵敏度决定了精度}：
        问题三 证明 $x(t) \to K_{\text{eff}}$ 恒成立——轨迹的归宿是确定的；
        问题四 揭示，如果 $K$ 估计偏差 $10\%$，那么``归宿''本身偏移 $\sim 814$ 吨/天
        （$\varepsilon_{K_{\text{eff}}, K} = 1.0000$，完全线性传导）。
        \textbf{系统轨迹将确定性地收敛至一个错误的吸引子}。在动力系统术语中，
        这不是 Lyapunov 稳定性的丧失，而是\textbf{吸引子位置的结构性偏移}——
        一种在传统稳定性分析框架下完全不可见但后果同样严重的失效模式。

  \item \textbf{收敛速率的安全裕度充足}：
        问题三 从性能约束出发定义了 $\beta$ 的运营上限：
        $\beta_{\max} = \gamma/r_{\min} - 1 = 3.40$。
        当前 $\beta = 1.59$ 距此边界尚有 $1.81$ 的裕度（需增大 $+114\%$ 才触及底线）。
        问题四 进一步验证了这一裕度的可靠性：即使 $\beta$ 增大 $10\%$（增至 $1.75$），
        $R^2$ 仅从 $0.9915$ 降至 $0.9882$，系统行为几乎不变。
        从动力系统视角看，系统距 $r = 0$ 的分岔点非常远（$r/r_{\min} \approx 1.70$），
        远离\textbf{临界慢化}（critical slowing down）区域——特征时间尺度
        $\tau \approx 118$ 天在参数扰动下保持稳定，不存在``越接近分岔点收敛越慢''
        的风险。\textbf{短期内不必担忧收敛过慢导致运营失效}。

  \item \textbf{真正的风险不在分岔失稳，而在吸引子错位}：
        问题三 的安全分析排除了系统从稳定变为不稳定的可能（不存在鞍-结分岔、
        跨临界分岔不能穿越 $r = 0$、叉形分岔因缺乏对称性而不适用）。
        但 问题四 揭示了一个更加现实的威胁：\textbf{基准容量 $K$ 是量级上的主导参数}——
        $\pm 10\%$ 偏差使 $R^2$ 从 $0.9915$ 急剧下降至 $\sim 0.3$，
        其 $R^2$ 降幅（$0.7285$）是次敏感参数 $\gamma$（$0.0105$）的约 70 倍、
        $\beta$（$0.0038$）的约 190 倍，$\alpha$ 的影响相比之下可忽略不计。在动力系统语言中，$K$ 直接决定了
        \textbf{唯一物理吸引子的位置}。当所有数据点密集分布在饱和平台附近
        （$x(365)/K_{\text{eff}} \approx 98.3\%$），吸引子位置的任何偏移
        都会被所有晚期数据点``集体感知''，产生系统性同向残差——这解释了
        $R^2$ 急剧下降的动力学机制。

        运营含义：\textbf{回收体系最大的风险不是系统失稳，而是对容量上限的
        误判导致的规划失误}——低估 $K$ 则规划产能不足、高估 $K$ 则造成
        资源浪费。这两种失败均可在拟合指标正常的表象下发生——所有常规统计检验均能通过。
\end{enumerate}

\subsubsection{关键结论与长效调控策略}

基于 问题三 稳定性分析和 问题四 灵敏度分析的联合发现，
为``沪尚回收''体系的长效稳定运行提出以下定量策略建议：

\paragraph*{1. 优先级一：精确测定基准容量 $K$（最高回报）}

$K$ 是模型的核心支柱——$\pm 10\%$ 偏差直接使 $R^2$ 从 $0.9915$ 急剧下降至 $\sim 0.3$，
敏感度远超其他参数（$\times 30$--$170$）。然而 $K$ 恰恰是四个参数中\textbf{最难直接测量}的：
它涉及常住人口总量、人均可回收物产生系数、低价值品类占比、回收网络地理覆盖率等外生变量，
任何一个估计环节的偏差都会沿 $\varepsilon(K_{\text{eff}}, K) = 1$ 的比例关系
直接传导到模型预测。

\textbf{建议措施}：
\begin{itemize}
  \item 开展全市可回收物潜力普查（抽样调查 + 遥感人口密度估算），
        将 $K$ 的估计误差控制在 $\pm 5\%$ 以内——这可将 $R^2$ 损失限制在可接受范围。
  \item 每年更新 $K$ 的估计值（因人口流动、消费结构变化等因素会改变基准容量），
        将模型重新校准纳入年度运维计划。
\end{itemize}

\paragraph*{2. 优先级二：精准降低抑制系数 $\beta$（最高效的调控杠杆）}

降低 $\beta$（改善运营效率）对恢复速度的影响是提升 $\alpha$（扩大促进）
对容量影响的 $\sim 38$ 倍（$0.614 / 0.0161 \approx 38$）。
且题目背景中明确指出了三大抑制源头：

\begin{itemize}
  \item \textbf{清运不及时}（智能回收箱满溢后居民无法投递）：
        优化清运调度算法，基于满溢传感器实现动态路线规划，将平均清运延迟控制在阈值以下。
  \item \textbf{老年群体操作困难}（智能设备适配性不足）：
        增设人工辅助回收点与智能设备并行；开发语音交互、大字体简化界面等适老化改造；
        在老龄化程度高的社区保留或增加传统回收渠道。
  \item \textbf{低价值品类回收意愿低}（源头分类不精细）：
        试行低价值可回收物补贴机制（按重量阶梯补贴到社区或居民），
        降低居民对低价值品类的分类心理成本。
\end{itemize}

当前 $\beta$ 已使有效增长率降至固有水平的 $38.6\%$，
说明抑制效应远大于促进效应——\textbf{运维阻力是当前回收体系的首要效率限制因素}。
然而 问题四 也提示：$\beta$ 弹性为 $-0.61$，边际收益递减，
应将资源集中在消除\textbf{最严重的瓶颈环节}
（如清运延迟和适老化），而非追求面面俱到的微优化。

\paragraph*{3. 优先级三：重新审视促进策略 $\alpha$ 的着力点（当前策略需转型）}

$\alpha$ 的弹性仅为 $0.016$——全市已建成 1.5 万个服务点 + 4300 台智能回收箱 + 800 个惠民服务点
后，站点数量增加对容量的边际贡献几乎为零。\textbf{但这并不意味着宣传引导没有价值}——
它可能通过以下渠道间接发挥作用：

\begin{itemize}
  \item \textbf{提升 $\gamma$}：宣传引导提高居民基础参与度，推升固有增长率，
        而非扩大饱和容量。应设计 AB 测试验证宣传投入对 $\gamma$ 的边际效应。
  \item \textbf{降低 $\beta$}：知识普及降低居民分类错误率（减少二次分拣的运维负担），
        社区动员提高居民对满溢和故障的主动上报率（间接改善清运效率）。
  \item \textbf{优化品类结构}：针对低价值品类的专项宣传可能改变 $K$ 的品类构成，
        而非 $K_{\text{eff}}$ 的总量。
\end{itemize}

\textbf{建议措施}：
\begin{itemize}
  \item 以``增加站点数量''为核心考核指标的扩张策略已进入边际收益递减区间——
        在 $K_{\text{eff}} \approx 8142$ 的数据硬约束下，$\alpha$ 对容量的弹性仅 $0.016$，
        继续扩大站点密度对总回收量的提升极为有限。
  \item 将促进工作的考核重心从``建了多少点''转向``每个点服务了多少人、
        实际投递频次和居民满意度''——关注促进策略对 $\gamma$ 和 $\beta$ 的间接效应，
        而非 $\alpha$ 本身。
  \item 将释放出的资源（人力、资金、管理精力）重新分配到 $\beta$ 的降低上
        （清运优化、适老化改造），实现\textbf{从规模扩张向质量提升的战略转型}。
\end{itemize}

\paragraph*{4. 优先级四：$\gamma$ 保持监测即可（低维护成本）}

$\gamma$ 的扰动影响中等（$R^2$ 保持在 $0.98$ 以上），弹性为 1（线性传递，无放大），
且 $\gamma$ 受制于居民回收习惯等慢变量，短期波动有限。
建议每季度通过小样本抽样调查验证 $\gamma$ 的稳定性，无需高频精确跟踪。

\paragraph*{5. 总体资源分配建议}

综合 问题三 和 问题四 的全部量化证据，建议``沪尚回收''体系的长效运维资源按以下优先级分配：

\begin{table}[htbp]
\centering
\caption{分层资源分配与长效运维策略}
\label{tab:resource_allocation}
\begin{tabular}{clp{7cm}}
\toprule
优先级 & 参数 & 行动方向与预期效果 \\
\midrule
\textbf{1} & $K$（基准容量） &
  \textbf{精确测量}：开展潜力普查，将估计误差压缩至 $\pm 5\%$ 内。
  预期效果：确保模型预测的可靠性基础不被摧毁。 \\

\textbf{2} & $\beta$（抑制） &
  \textbf{定向削减}：优化清运调度 + 适老化改造 + 低价值品类补贴。
  预期效果：每降低 $\beta$ 约 $10\%$，有效增长率 $r$ 提升约 $6.5\%$。 \\

\textbf{3} & $\alpha$（促进） &
  \textbf{战略转型}：从扩大站点数量转向提升站点服务质量，
  考核指标从覆盖率转向人均使用频次和满意度。关注促进策略对 $\gamma$、$\beta$ 的间接效应。
  预期效果：释放无效扩张的沉没成本，将资源重新导向高回报领域。 \\

\textbf{4} & $\gamma$（固有增长率） &
  \textbf{常规监测}：季度抽样验证，无需精确跟踪。系统对 $\gamma$ 的偏差具有鲁棒性。
  预期效果：维持对居民参与度长期趋势的基本感知。 \\
\bottomrule
\end{tabular}
\end{table}

\section{模型评价与改进}

\subsection{模型优点}

\begin{enumerate}
    \item \textbf{建模思路递进清晰}：从仅含容量约束的基础 Logistic 模型（问题一），到引入拮抗调节参数的扩展模型（问题二），再到动力系统稳定性分析（问题三）和参数灵敏度量化（问题四），四步递进完整覆盖了``建模—拟合—稳定性判定—灵敏度控制''的全流程，每步均建立在前一步的发现之上。
    \item \textbf{参数物理意义明确}：$\gamma$（固有增长率）对应居民基础参与度，$\alpha$（促进系数）对应站点扩张与宣传引导的容量放大效应，$\beta$（抑制系数）对应清运延迟与运维不足的速率压缩效应，$K$（基准容量）对应常住人口与人均可回收物产生量——每个参数均可与具体管理措施直接对应，使模型的定量结论具有可操作性。
    \item \textbf{结论有严格的数学保证，且提供了独特的风险洞察}：区别于纯统计方法，本文关于系统安全性的每一条结论——$\lambda = -r < 0$ 的渐近稳定性、跨临界分岔的不可穿越性、鞍-结分岔与叉形分岔的不适用性——均可追溯到严格的数学论证，而非数值实验的经验推断。在此基础上，动力系统框架进一步揭示了``吸引子错位''（$K$ 估计偏差）这一传统统计诊断无法发现的核心风险——系统会以可靠的收敛性趋向一个错误的目标值，这是在分岔分析证明系统结构稳定之后暴露出来的更深层的脆弱性。
\end{enumerate}

\subsection{模型不足与改进方向}

\begin{enumerate}
    \item \textbf{时不变假设}：模型将促进系数 $\alpha$ 和抑制系数 $\beta$ 视为阶段内恒定，未刻画其随时间演化的动态过程。改进方向：可将 $\alpha(t)$ 和 $\beta(t)$ 建模为时变函数（如 logistic 衰减型或分段常数），引入滑动窗口估计等系统辨识方法跟踪参数的时变轨迹。

    \item \textbf{品类同质化假设}：模型以总回收量 $x(t)$ 为单一状态变量，未区分纸类、塑料、金属、玻璃等品类的差异化回收规律。改进方向：可构建多品类耦合 ODE 系统 $\dot{x}_i = f_i(x_1, \dots, x_n)$，各品类的 $\gamma_i$、$K_i$ 独立估计，并通过交互项刻画品类间的替代或互补效应。

    \item \textbf{空间均质化假设}：模型基于全市总量数据，未体现中心城区与郊区、不同街道间的空间异质性。改进方向：可引入分区耦合的 ODE 系统 $\dot{x}_j = f_j(x_j) + \sum_{k \neq j} D_{jk}(x_k - x_j)$，分析区域间回收量的扩散与竞争效应，为站点布局与清运路线优化提供支持。

    \item \textbf{数据量有限}：模型基于 9 个数据点拟合，参数置信区间相对较宽（已通过残差 Bootstrap 验证其稳健性）。改进方向：引入贝叶斯更新框架逐步缩小参数不确定度，或采用 Sobol' 全局灵敏度分析替代 OAT 局部方法，量化参数间交互效应对预测方差的贡献。
\end{enumerate}
\subsection{模型适用性说明}

本文模型的核心假设——一维自治 Logistic 动力学——适用于以下场景：（1）回收体系已建立并持续运营，回收量呈单调增长趋势；（2）体系的总容量在一定时期内保持稳定；（3）促进与抑制因素的相对强度在分析周期内不发生根本性逆转。对于体系初创期（回收量可能波动或下降）、政策剧烈调整期（参数可能发生结构性突变）、或需要在品类/空间维度上进行精细化分析的场景，建议采用上述改进方向中的扩展模型。

\section{结论}

本文围绕上海``沪尚回收''体系回收量的动态演化规律，构建了递进式 Logistic 动力学模型，完成从基础拟合到系统安全性证明、再到参数灵敏度量化的完整分析链条。核心发现如下。

\textbf{1. 抑制效应是当前体系的首要效率限制因素。}
$\beta = 1.59$ 将有效增长率压缩至固有水平的 $38.6\%$，而 $\alpha = 0.0164$ 对饱和容量的提升仅有 $1.64\%$——清运延迟、设备适老化不足等运维阻力对系统效率的拖累远大于站点扩张的边际贡献。在全市 1.5 万个服务点与 4300 台智能回收箱的覆盖密度下，继续扩站的边际收益已趋近于零。

\textbf{2. 系统不发生分岔失稳，但存在更隐蔽的吸引子偏移风险。}
问题三 逐一排除了鞍-结、跨临界（物理锁定于 $r>0$）、叉形分岔三种失稳机制，唯一物理平衡点恒为渐近稳定。然而问题四 揭示：$K \pm 10\%$ 即可使 $R^2$ 从 $0.9915$ 骤降至 $\sim 0.3$——系统会以可靠的收敛性趋向一个错误的吸引子，这是传统稳定性检验完全不可见的失效模式。

\textbf{3. 参数灵敏度呈现极化的``$K$-主导''格局。}
$K$ 的 $R^2$ 降幅（$0.7285$）是 $\gamma$（$0.0105$）的约 70 倍、$\beta$（$0.0038$）的约 190 倍，$\alpha$ 的影响可忽略。所有解析弹性绝对值均 $\leq 1$，模型在局部意义下良态。$\beta$ 弹性 $|\varepsilon_{r,\beta}| = 0.61$ 提示抑制效应的边际收益递减——降低 $\beta$ 有效但不宜期待单点突破。

\textbf{4. 长效运营应遵循分层资源分配策略。}
优先级依次为：（1）精确普查 $K$，将估计误差压缩至 $\pm 5\%$ 内——这是模型预测可靠性的前提；（2）定向降低 $\beta$，优化清运调度与适老化改造，每降低 $\beta$ 约 $10\%$ 可使 $r$ 回升约 $6.5\%$；（3）促进策略从规模扩张转向质量提升，关注宣传引导对 $\gamma$ 和 $\beta$ 的间接效应；（4）$\gamma$ 季度抽样监测即可，系统对其偏差具有充分鲁棒性。

\textbf{5. 局限与展望。}
本文未刻画 $\alpha$、$\beta$ 的时变特性与区域/品类间的异质性。后续可沿三个方向拓展：将参数推广为时变函数以捕捉政策干预的阶段性效果；引入分区耦合 ODE 系统分析区域间扩散效应；按品类拆分模型识别差异化增长规律。

\newpage
\begin{thebibliography}{99}

\bibitem{problem} 2026 年数学建模训练赛 A 题及附件数据[Z]. 2026.

\bibitem{openai_codex} OpenAI Codex, GPT-5, OpenAI, 2026-07-30.

\bibitem{verhulst} VERHULST P F. Notice sur la loi que la population suit dans son accroissement[J]. Correspondance Math\'{e}matique et Physique, 1838, 10: 113--121.

\bibitem{efron} EFRON B, TIBSHIRANI R J. An Introduction to the Bootstrap[M]. New York: Chapman \& Hall, 1993.

\bibitem{strogatz} STROGATZ S H. Nonlinear Dynamics and Chaos: With Applications to Physics, Biology, Chemistry, and Engineering[M]. 3rd ed. Boca Raton: Chapman \& Hall/CRC, 2024.

\bibitem{saltelli} SALTELLI A, RATTO M, ANDRES T, et al. Global Sensitivity Analysis: The Primer[M]. Chichester: John Wiley \& Sons, 2008. DOI: 10.1002/9780470725184.

\end{thebibliography}

\clearpage
\appendix

\section{支撑材料文件清单与复现说明}

\subsection{支撑材料文件清单}

支撑材料压缩包按“代码、结果图、AI 使用详情”三类组织。赛题提供的 9 个观测点及固定参数
已逐项写入程序，未另行重复提交赛题原始附件；本文未使用自主查阅的外部数据集。

\begin{table}[H]
\centering
\small
\begin{tabular}{p{4.6cm}p{2.7cm}p{5.9cm}}
\toprule
文件 & 对应内容 & 作用与主要输出 \\
\midrule
\path{Q1/logistic_model.py} & 问题一 & 基础 Logistic 预测、误差指标与 \path{logistic_fit.png} \\
\path{Q2/advanced_model.py} & 问题二 & 约束非线性拟合、Delta 方法与 \path{model_q2.png} \\
\path{Q2/bootstrap_validation.py} & 问题二 & $B=10\,000$ 次残差 Bootstrap 交叉验证 \\
\path{Q3/stability_analysis.py} & 问题三 & 平衡点、特征值、性能边界与 \path{stability_analysis.png} \\
\path{Q4/sensitivity_analysis.py} & 问题四 & $\pm10\%$ OAT 分析及两幅灵敏度结果图 \\
\path{requirements.txt} & 运行环境 & Python 第三方依赖及最低版本 \\
\path{README.md} & 复现说明 & 运行顺序、命令、输入与输出说明 \\
{\codefont AI工具使用详情.pdf} & AI 使用记录 & 工具、用途、关键交互及人工修改情况 \\
\bottomrule
\end{tabular}
\caption{支撑材料文件列表}
\label{tab:supporting-files}
\end{table}

\subsection{运行环境与复现顺序}

建议使用 Python 3.10 或更高版本，安装 NumPy、SciPy 与 Matplotlib 后，依次进入
\texttt{Q1}、\texttt{Q2}、\texttt{Q3}、\texttt{Q4} 目录运行表
\ref{tab:supporting-files} 中的脚本。各脚本均内置赛题数据与固定随机种子，无需联网，
可独立运行；图像文件输出到脚本所在目录。推荐运行顺序为问题一、问题二（先
\texttt{advanced\_model.py}，后 \texttt{bootstrap\_validation.py}）、问题三、问题四。

\section{AI工具使用说明}

本文使用 OpenAI Codex（GPT-5）辅助核对参考文献著录信息、检查 LaTeX 引用闭环、
根据竞赛格式规范完善附录、复核程序可运行性，并对相关说明文字进行规范化。
AI 未代替参赛队作出模型选择，未更改观测数据或数值结果，也未替代参赛队对推导和结论
进行验证。参赛队对最终论文的原创性、真实性和准确性承担全部责任。工具版本、使用日期、
关键提示词与回复、采纳及人工修改情况，另见支撑材料中的“AI工具使用详情.pdf”。
如在提交前继续使用其他 AI 工具或开展新的关键交互，须同步补充该文件及参考文献记录。

\section{关键Python代码摘录}

\subsection{问题一：基础 Logistic 模型}

以下仅保留解析模型与误差评价的核心实现；绘图、控制台输出和报告生成代码从略。
\lstinputlisting[language=Python,firstline=23,lastline=41,firstnumber=23]{../Q1/logistic_model.py}

\subsection{问题二：改进模型与参数拟合}

以下代码完成带物理约束的非线性拟合、参数反解和 Delta 方法误差传播。
\lstinputlisting[language=Python,firstline=47,lastline=78,firstnumber=47]{../Q2/advanced_model.py}

\subsection{问题二：残差 Bootstrap 验证}

以下为固定随机种子、残差重抽样、重复拟合和经验区间计算的关键循环。
\lstinputlisting[language=Python,firstline=68,lastline=117,firstnumber=68]{../Q2/bootstrap_validation.py}

\subsection{问题三：平衡点与稳定性分析}

以下代码给出平衡点特征值、松弛时间和运营安全边界的计算。
\lstinputlisting[language=Python,firstline=66,lastline=98,firstnumber=66]{../Q3/stability_analysis.py}

\subsection{问题四：OAT 参数灵敏度分析}

以下为四参数 $\pm10\%$ 单因素扰动、评价指标计算与解析弹性计算。
\lstinputlisting[language=Python,firstline=55,lastline=114,firstnumber=55]{../Q4/sensitivity_analysis.py}

为控制论文附录篇幅，以上省略了重复的数据定义、绘图和结果文件写出语句。
五个完整、可直接运行的原始程序均收录于独立提交的支撑材料压缩包中，
文件路径与运行顺序见附录 A。


\end{document}
